\documentclass{book}
\usepackage[backend=biber,natbib=true,style=authoryear]{biblatex}
\addbibresource{/home/nguyen/1_NQBH/math/bib.bib}
\usepackage[utf8]{inputenc}
\usepackage{float}
\usepackage{graphicx}
\usepackage[colorlinks=true,linkcolor=blue,urlcolor=red,citecolor=magenta]{hyperref}
\usepackage{amsmath,amssymb,amsthm}
\allowdisplaybreaks
\numberwithin{equation}{section}
\newtheorem{assumption}{Assumption}[section]
\newtheorem{lemma}{Lemma}[section]
\newtheorem{corollary}{Corollary}[section]
\newtheorem{definition}{Definition}[section]
\newtheorem{proposition}{Proposition}[section]
\newtheorem{theorem}{Theorem}[section]
\newtheorem{notation}{Notation}[section]
\newtheorem{remark}{Remark}[section]
\newtheorem{example}{Example}[section]
\newtheorem{ques}{Question}[section]
\newtheorem{problem}{Problem}[section]
\newtheorem{conjecture}{Conjecture}[section]
\usepackage[left=0.5in,right=0.5in,top=1.5cm,bottom=1.5cm]{geometry}
\usepackage{fancyhdr}
\pagestyle{fancy}
\fancyhf{}
\addtolength{\headheight}{0pt} % obsolete
\lhead{\itshape \scriptsize \chaptername~\thechapter}
\rhead{\itshape \scriptsize \nouppercase{\leftmark}} %\nouppercase !
\renewcommand{\chaptermark}[1]{\markboth{#1}{}}
\cfoot{\thepage}

\title{Around the World}
\author{Collector: Nguyen Quan Ba Hong}
\date{\today}

\begin{document}
\maketitle
\setcounter{secnumdepth}{6}
\setcounter{tocdepth}{6}
\tableofcontents

%------------------------------------------------------------------------------%

\part{Austria}

\chapter{\href{https://en.wikipedia.org/wiki/Austria}{Wikipedia/Austria}}


%------------------------------------------------------------------------------%

\paragraph{France}

\chapter{\href{https://en.wikipedia.org/wiki/France}{Wikipedia/France}}



%------------------------------------------------------------------------------%

\part{Germany}

\chapter{\href{https://en.wikipedia.org/wiki/Germany}{Wikipedia/Germany}}

%------------------------------------------------------------------------------%

\chapter{\href{https://en.wikipedia.org/wiki/Weierstrass_Institute}{Wikipedia/Weierstrass Institute}}

\textbf{Weierstrass Institute.}
\begin{itemize}
    \item \textbf{Abbreviation.} WIAS
    \item \textbf{Predecessor.} Karl Weierstrass Institute for Mathematics of the GDR's Academy of Science.
    \item \textbf{Formation.} 1992.
    \item \textbf{Type.} Research Institute.
    \item \textbf{Purpose.} Pure and applied mathematical research.
    \item \textbf{Location.} Berlin.
    \item \textbf{Director.} Michael Hintermüller
    \item \textbf{Parent organization.} Forschungsverbund Berlin e.V.
    \item \textbf{Affiliations.} \href{https://en.wikipedia.org/wiki/Leibniz_Association}{Leibniz Association}.
    \item \textbf{Staff.} $\sim 150$.
    \item \textbf{Website.} \href{http://www.wias-berlin.de/index.jsp?lang=1}{www.wias-berlin.de}
\end{itemize}
The \textit{Weierstrass Institute}, formally the \textit{Weierstrass Institute for Applied Analysis and Stochastics} (WIAS), is a part of the Forschungsverbund Berlin e.V.[``Weierstrass Institute for Applied Analysis and Stochastics (WIAS), Leibniz Institute in Forschungsverbund Berlin e.V. - Forschungsverbund Berlin''. www.fv-berlin.de. Retrieved Feb 25, 2016.] and a member of the \href{https://en.wikipedia.org/wiki/Leibniz_Association}{Leibniz Association}.[``Leibniz Gemeinschaft: Institutes \& Museums/WIAS''. \url{www.leibniz-gemeinschaft.de}. Retrieved Feb 25, 2016.]

Based in Berlin's district Mitte, the institute's research activities involve \href{https://en.wikipedia.org/wiki/Applied_mathematics}{applied} and \href{https://en.wikipedia.org/wiki/Pure_mathematics}{pure mathematics}.

%
Since Feb 2011, the \href{https://en.wikipedia.org/wiki/International_Mathematical_Union}{International Mathematical Union} (IMU) Secretariat has been located in Berlin at WIAS.[``International Mathematical Union (IMU): Secretariat''. \url{www.mathunion.org}. Retrieved Feb 25, 2016.]

\section{History}
The institute developed from the ``Karl Weierstrass Institute for Mathematics'' of the former \href{https://en.wikipedia.org/wiki/GDR}{GDR}'s \href{https://en.wikipedia.org/wiki/German_Academy_of_Sciences_at_Berlin}{Academy of Sciences at Berlin}.

Following a recommendation of the \href{https://en.wikipedia.org/wiki/German_Council_of_Science_and_Humanities}{German Council of Science and Humanities}, the institute was established on Jan 1, 1992.

%
The institute is named after the Berlin mathematician \href{https://en.wikipedia.org/wiki/Karl_Weierstra%C3%9F}{Karl Weierstraß} (1815-1897), who was well known for his profound processing of analysis.

\section{Tasks}
The main task of WIAS is the implementation of project-oriented research in \href{https://en.wikipedia.org/wiki/Applied_mathematics}{applied mathematics}, especially in applied \href{https://en.wikipedia.org/wiki/Analysis}{analysis} and \href{https://en.wikipedia.org/wiki/Stochastics}{stochastics}.

The research activity is oriented on specific application situations and is confirmed by cooperations with scientific institutions and the economy.

It includes the entire spectrum of the solution, from \href{https://en.wikipedia.org/wiki/Mathematical_model}{mathematical modeling} to mathematical-theoretical analysis of models, up to the development of \href{https://en.wikipedia.org/wiki/Algorithm}{algorithms} and numerical simulation of technological processes.

%
Research at WIAS focuses on the following application areas:
\begin{itemize}
    \item \href{https://en.wikipedia.org/wiki/Nanoelectronics}{Nano}- and \href{https://en.wikipedia.org/wiki/Optoelectronics}{Optoelectronics},
    \item \href{https://en.wikipedia.org/wiki/Mathematical_optimization}{Optimization} and \href{https://en.wikipedia.org/wiki/Control_theory}{control} in \href{https://en.wikipedia.org/wiki/Technology}{technology} and economy,
    \item Materials modeling,
    \item Flow and transport,
    \item \href{https://en.wikipedia.org/wiki/Energy_transformation}{Conversion}, \href{https://en.wikipedia.org/wiki/Energy_storage}{storage} and \href{https://en.wikipedia.org/wiki/Electric_power_distribution}{distribution} of \href{https://en.wikipedia.org/wiki/Energy}{energy},
    \item Random phenomena in \href{https://en.wikipedia.org/wiki/Nature}{nature} and \href{https://en.wikipedia.org/wiki/Economy}{economy}.
    
    [``Leibniz Gemeinschaft: Institutes \& Museums/WIAS''. \url{www.leibniz-gemeinschaft.de}. Retrieved Feb 25, 2016.]
\end{itemize}
Within these application areas, questions are investigated, that play a central role in the development of key technologies (e.g. \href{https://en.wikipedia.org/wiki/Materials_science}{materials science}, \href{https://en.wikipedia.org/wiki/Manufacturing_engineering}{manufacturing engineering}, \href{https://en.wikipedia.org/wiki/Biomedical_engineering}{biomedical engineering}), as well as applications in the economy (e.g. \href{https://en.wikipedia.org/wiki/Finance}{finances}).

%
Services are not part of the tasks.

\section{Cooperation}
WIAS cooperates closely with university and non-university institutions and is involved in numerous joint research projects, for which additional funds are raised within the framework of competition procedures (e. g. of the \href{https://en.wikipedia.org/wiki/German_Research_Foundation}{German Research Foundation} or the \href{https://en.wikipedia.org/wiki/European_Commission}{European Commission}.)[4][5][6]

%
Close relationships are maintained with the 3 Berlin universities, \href{https://en.wikipedia.org/wiki/Humboldt_University_of_Berlin}{Humboldt University}, \href{https://en.wikipedia.org/wiki/Technical_University_of_Berlin}{Technical University} and \href{https://en.wikipedia.org/wiki/Free_University_of_Berlin}{Free University}, through cooperation contracts and 6 joint appeals based on these contracts.

Together with the mathematical institutes of these universities and the \href{https://en.wikipedia.org/wiki/Zuse_Institute_Berlin}{Zuse Institute Berlin} WIAS operates the Berlin mathematical center of excellence, MATHEON.[Matheon. ``\textit{MATHEON Research Center}''. \url{matheon.de}. Retrieved Feb 26, 2016.]

%
Within the \href{https://en.wikipedia.org/wiki/Leibniz_Association}{Leibniz Association} WIAS coordinates a network ``Mathematical Modeling and Simulation''. [``Leibniz Gemeinschaft: Forschung/Leibniz-Netzwerke/Leibniz-Netzwerk ``\textit{Mathematische Modellierung und Simulation} (MMS)''''. \url{www.leibniz-gemeinschaft.de}. Retrieved Feb 25, 2016.]

\section{Infrastructure}
The total budget of the institute was 12.9 million euros in 2015, of which the federal and state governments, within the framework of the basic funding, financed 9.5 million euros equally. 

Approximately 150 people work at WIAS.[``About WIAS - Facts \& Figures''. \url{www.wias-berlin.de}. Retrieved Feb 25, 2016.]

\section{External links}
\begin{itemize}
    \item \href{http://www.wias-berlin.de/index.jsp?lang=1}{WIAS Berlin}
    \item \href{http://www.fv-berlin.de/}{Website of Forschungsverbund Berlin e.V.}\hfill$\square$
\end{itemize}

%------------------------------------------------------------------------------%

\part{Vietnam}

\chapter{\href{https://en.wikipedia.org/wiki/Vietnam}{Wikipedia/Vietnam}}



%------------------------------------------------------------------------------%

\printbibliography[heading=bibintoc]
\end{document}